%============================================================
\section{Conclusion}\label{sec:Conclusion}
%============================================================
% Previous conclusion:
% This paper introduces search as a new method for smoothing out discrete-continuous
% choice (DC) problems. While search can significantly reduce the size of the disconti-
% nuities in the decision rule of a DC problem, to get uniform smoothness around the
% discontinuities, we need to incorporate taste shocks as well. Together, search and taste
% shocks offer modelers greater control over the degree of smoothness in DC models, while
% rationalizing the smoothness into the agent's choice problem. With ample smoothness,
% discontinuities in the decision rule can be smoothed out entirely; thus, eliminating the
% need for incorporating an additional, computationally costly, step of identifying and
% removing suboptimal points generated in the EGM step for DC models. This paper
% studies search in the context of a consumption-saving-retirement problem; however, the
% framework can be easily applied to models of housing purchase, lumpy firm investments,
% and a wide array of other DC models.
% We additionally document how a coarse income-shock discretization injects phantom
% kinks into the decision rule, and describe a transition-matrix method that integrates
% the continuous income density and removes those numerical artifacts at moderate cost.
Directed search and a rejectable price draw can match the smoothness that Extreme
Value Type~I taste shocks give to a discrete-continuous problem. The agent chooses
effort in the unit interval at a convex utility cost. Effort raises the probability that
an option to switch arrives. The draw is the content of that option: an account factor
in the retirement model, and the purchase price of the new stock in the durable model.
A bad draw is rejected. Search alone leaves one side of the threshold unsmoothed,
because effort is zero when the switch is not strictly preferred. The price draw fills
that side. The option value of the draw raises effort.

The retirement model illustrates this mechanism. The quantitative comparison is the
keep-versus-adjust model of durables. Taste shocks, price dispersion, and search all
smooth the $(S,s)$ band. In steady state they move the adjustment probability in
opposite directions. A larger taste scale raises it, and durables and assets fall.
Wider price dispersion raises it, because better draws are accepted; durables rise and
assets fall. A higher search cost lowers it; durables fall and assets rise. Search and
price dispersion can therefore be disciplined separately. Impulse responses of
consumption and of the adjustment probability, at a common adjustment frequency,
are the open exercise.

Search acts near the $(S,s)$ thresholds, so the Euler equation still describes agents
far from indifference. Optimal effort maps into an expected search duration. A search
cost specified to scale with permanent income does not add a state to the normalized
problem. An Extreme Value shock of the same scale for every household does not scale
in that way.
